\documentclass{article}
\usepackage{amsmath,amssymb,amsthm}
\usepackage{algorithm}
\usepackage{algpseudocode}
\usepackage{enumitem}
\usepackage{hyperref}
\usepackage{geometry}
\geometry{margin=1in}
\newtheorem{theorem}{Theorem}
\newtheorem{lemma}{Lemma}
\newtheorem{assumption}{Assumption}
\newtheorem{corollary}{Corollary}
\newtheorem{prop}{Proposition}
\newcommand{\prox}{\operatorname{prox}}
\newcommand{\R}{\mathbb{R}}

\begin{document}

\section*{Algorithm Name}
\textbf{Proximal Gradient with Momentum (PGM)}  

The method solves the composite convex problem  

\[
\min_{x\in\R^{d}} \; F(x)\;:=\; f(x)+g(x),
\]

where $f$ is convex and $L$‑smooth, $g$ is convex (possibly non‑smooth) and admits an easy proximal operator.

\section*{Mathematical Formulation}
Let $\eta>0$ be a stepsize and $\beta\in[0,1)$ a momentum parameter.  
Initialize $x_{0}\in\R^{d}$ and $v_{-1}=0$. For $k=0,1,2,\dots$ compute  

\[
\begin{aligned}
v_{k}      &:= \beta\,v_{k-1} + \nabla f(x_{k}) , \tag{1}\\[4pt]
x_{k+1}    &:= \prox_{\eta g}\bigl(x_{k}-\eta\, v_{k}\bigr). \tag{2}
\end{aligned}
\]

When $g\equiv0$, $\prox_{\eta g}$ reduces to the identity and (2) becomes a plain gradient step with a momentum‑augmented gradient.

\section*{Pseudocode}
\begin{algorithm}[H]
\caption{Proximal Gradient with Momentum (PGM)}
\begin{algorithmic}[1]
\Require Stepsize $\eta>0$, momentum $\beta\in[0,1)$, initial point $x_{0}\in\R^{d}$
\State $v_{-1}\gets 0$
\For{$k=0,1,2,\dots$}
    \State $v_{k}\gets \beta\,v_{k-1}+ \nabla f(x_{k})$ \Comment{momentum update}
    \State $x_{k+1}\gets \prox_{\eta g}\!\bigl(x_{k}-\eta\,v_{k}\bigr)$ \Comment{proximal step}
\EndFor
\end{algorithmic}
\end{algorithm}

\section*{Dry Run Example}
Consider the scalar problem  

\[
f(w)=(w-3)^{2},\qquad g\equiv0,
\]

so $\nabla f(w)=2(w-3)$.  
Take $\eta=0.1$, $\beta=0.5$, and $w_{0}=0$.  
Since $g\equiv0$, $\prox_{\eta g}(z)=z$.

\begin{enumerate}[label=Iteration \arabic*:, leftmargin=*]
\item $k=0$
\[
\begin{aligned}
v_{0}&=0.5\cdot0+2(w_{0}-3)=2(0-3)=-6,\\
w_{1}&=w_{0}-\eta v_{0}=0-0.1(-6)=0.6,\\
F(w_{1})&=(0.6-3)^{2}=5.76.
\end{aligned}
\]

\item $k=1$
\[
\begin{aligned}
v_{1}&=0.5\,v_{0}+2(w_{1}-3)=0.5(-6)+2(0.6-3)=-3-4.8=-7.8,\\
w_{2}&=w_{1}-0.1\,v_{1}=0.6-0.1(-7.8)=1.38,\\
F(w_{2})&=(1.38-3)^{2}=2.6244.
\end{aligned}
\]

\item $k=2$
\[
\begin{aligned}
v_{2}&=0.5\,v_{1}+2(w_{2}-3)=0.5(-7.8)+2(1.38-3)=-3.9-3.24=-7.14,\\
w_{3}&=w_{2}-0.1\,v_{2}=1.38-0.1(-7.14)=2.094,\\
F(w_{3})&=(2.094-3)^{2}=0.821.
\end{aligned}
\]
\end{enumerate}
The iterates move toward the minimizer $w^{*}=3$ and the objective values decrease.

\newpage
\section*{Assumptions}
\begin{assumption}[Convexity and Smoothness of $f$]\label{as:convexsmooth}
$f:\R^{d}\to\R$ is convex and differentiable with $L$‑Lipschitz gradient, i.e.  

\[
\|\nabla f(x)-\nabla f(y)\|\le L\|x-y\|,\qquad\forall x,y\in\R^{d}.
\] 
\end{assumption}

\begin{assumption}[Convexity of $g$]\label{as:gconvex}
$g:\R^{d}\to\R\cup\{+\infty\}$ is proper, closed and convex. Its proximal operator  

\[
\prox_{\eta g}(z):=\arg\min_{x}\Bigl\{ g(x)+\frac{1}{2\eta}\|x-z\|^{2}\Bigr\}
\]

is assumed to be efficiently computable.
\end{assumption}

\begin{assumption}[Step size and momentum]\label{as:params}
The stepsize satisfies  

\[
0<\eta\le \frac{1}{L},
\]

and the momentum parameter obeys  

\[
0\le\beta<1.
\]
\end{assumption}

\section*{Main Convergence Theorem}
\begin{theorem}[O(1/k) function value rate]\label{thm:rate}
Let $\{x_{k}\}$ be generated by the PGM algorithm under Assumptions \ref{as:convexsmooth}–\ref{as:params}.  
Denote $F^{*}:=\inf_{x}F(x)$ and let $x^{*}\in\arg\min F$. Then for all $K\ge 1$

\[
F(x_{K})-F^{*}\;\le\;\frac{\|x_{0}-x^{*}\|^{2}}{2\eta\,(1-\beta)\,K}.
\tag{3}
\]

Consequently $\displaystyle\lim_{k\to\infty}F(x_{k})=F^{*}$ and the rate is $O(1/k)$.
\end{theorem}

\section*{Theoretical Properties}
\begin{enumerate}[label=\arabic*.]
\item \textbf{Stability.} The bound (3) shows that the iterates remain bounded as long as the initial distance $\|x_{0}-x^{*}\|$ is finite.
\item \textbf{Hyper‑parameter sensitivity.} The factor $1/(1-\beta)$ appears in the rate; larger momentum (i.e., $\beta\!\uparrow\!1$) slows the theoretical guarantee, although in practice a moderate $\beta\in[0.5,0.9]$ often accelerates convergence.
\item \textbf{Comparison to baselines.}
\begin{itemize}
\item With $\beta=0$ the method reduces to the classical proximal gradient algorithm, and (3) recovers the standard $O(1/k)$ bound.
\item For $\beta>0$ the algorithm enjoys the same asymptotic rate but may exhibit smaller constants in practice because the momentum term reduces oscillations.
\end{itemize}
\item \textbf{Special case $g\equiv0$.} When $g$ is absent, (3) simplifies to the bound for the (momentum‑augmented) gradient method.
\end{enumerate}

\section*{Preliminaries (Definitions and Lemmas)}
\begin{lemma}[Descent Lemma]\label{lem:descent}
If $f$ has $L$‑Lipschitz gradient, then for any $x,y\in\R^{d}$  

\[
f(y)\le f(x)+\langle\nabla f(x),y-x\rangle+\frac{L}{2}\|y-x\|^{2}.
\tag{4}
\]
\end{lemma}
\begin{proof}
Standard consequence of the integral form of the gradient; see e.g. \cite{Nesterov2004}.
\end{proof}

\begin{lemma}[Proximal Inequality]\label{lem:prox}
For any $z\in\R^{d}$ and any $x$, the proximal operator satisfies  

\[
g(\prox_{\eta g}(z))+\frac{1}{2\eta}\bigl\|\prox_{\eta g}(z)-z\bigr\|^{2}
\le g(x)+\frac{1}{2\eta}\|x-z\|^{2}
-\frac{1}{2\eta}\bigl\|x-\prox_{\eta g}(z)\bigr\|^{2}.
\tag{5}
\]
\end{lemma}
\begin{proof}
Follows from the optimality condition of the minimisation defining $\prox_{\eta g}(z)$.
\end{proof}

\section*{Main Proof (Step‑by‑Step Derivation)}
We prove Theorem \ref{thm:rate} by tracking the squared distance to an optimal point $x^{*}$.

\noindent\textbf{Notation.} Throughout let $x^{*}\in\arg\min F$ and define  

\[
\Delta_{k}:=\frac{1}{2}\|x_{k}-x^{*}\|^{2}.
\]

\noindent\textbf{Step 1.} Expand $\Delta_{k+1}$ using the update (2).  

\[
\begin{aligned}
\Delta_{k+1}
&=\frac12\bigl\|\,\prox_{\eta g}(x_{k}-\eta v_{k})-x^{*}\bigr\|^{2}\\
&\overset{[Step 1]}{\le}
\frac12\Bigl\|x_{k}-\eta v_{k}-x^{*}\Bigr\|^{2}
       -\frac12\bigl\|\,\prox_{\eta g}(x_{k}-\eta v_{k})-(x_{k}-\eta v_{k})\bigr\|^{2}  \\
&\qquad\text{(by Lemma \ref{lem:prox} with $z=x_{k}-\eta v_{k}$ and $x=x^{*}$)}\\
&=\Delta_{k}-\eta\langle v_{k},x_{k}-x^{*}\rangle
   +\frac{\eta^{2}}{2}\|v_{k}\|^{2}
   -\frac12\bigl\|\,\prox_{\eta g}(x_{k}-\eta v_{k})-(x_{k}-\eta v_{k})\bigr\|^{2}.
\tag{6}
\end{aligned}
\]

\noindent\textbf{Step 2.} Replace $v_{k}$ by its definition (1).  

\[
\begin{aligned}
\langle v_{k},x_{k}-x^{*}\rangle
&=\langle\beta v_{k-1}+\nabla f(x_{k}),\,x_{k}-x^{*}\rangle\\
&=\beta\langle v_{k-1},x_{k}-x^{*}\rangle+\langle\nabla f(x_{k}),x_{k}-x^{*}\rangle .
\tag{7}
\end{aligned}
\]

\noindent\textbf{Step 3.} Use convexity of $f$ to bound the second term in (7).  

\[
\langle\nabla f(x_{k}),x_{k}-x^{*}\rangle
\ge f(x_{k})-f(x^{*}) .
\tag{8}
\]

\noindent\textbf{Step 4.} Upper‑bound $\|v_{k}\|^{2}$.
From (1) and the Cauchy–Schwarz inequality,

\[
\begin{aligned}
\|v_{k}\|^{2}
&=\|\beta v_{k-1}+\nabla f(x_{k})\|^{2}\\
&\le (1+\beta)\bigl(\beta\|v_{k-1}\|^{2}+\|\nabla f(x_{k})\|^{2}\bigr)
\le \frac{1+\beta}{1-\beta}\,\|\nabla f(x_{k})\|^{2}
\tag{9}
\end{aligned}
\]

where we used the recursion $\|v_{k-1}\|\le \frac{1}{1-\beta}\|\nabla f(x_{k-1})\|$ that follows by unrolling (1) and the geometric series $\sum_{i=0}^{\infty}\beta^{i}=1/(1-\beta)$.

\noindent\textbf{Step 5.} Substitute (7)–(9) into (6).  

\[
\begin{aligned}
\Delta_{k+1}
&\le \Delta_{k}
      -\eta\Bigl[\beta\langle v_{k-1},x_{k}-x^{*}\rangle
                +f(x_{k})-f(x^{*})\Bigr] \\
&\qquad +\frac{\eta^{2}}{2}\,\frac{1+\beta}{1-\beta}\,\|\nabla f(x_{k})\|^{2}
      -\frac12\bigl\|\,\prox_{\eta g}(x_{k}-\eta v_{k})-(x_{k}-\eta v_{k})\bigr\|^{2}.
\tag{10}
\end{aligned}
\]

\noindent\textbf{Step 6.} Apply the descent lemma (Lemma \ref{lem:descent}) with $y=\prox_{\eta g}(x_{k}-\eta v_{k})$:

\[
\begin{aligned}
f\bigl(\prox_{\eta g}(x_{k}-\eta v_{k})\bigr)
&\le f(x_{k})+\langle\nabla f(x_{k}),\,
            \prox_{\eta g}(x_{k}-\eta v_{k})-x_{k}\rangle
      +\frac{L}{2}\bigl\|\prox_{\eta g}(x_{k}-\eta v_{k})-x_{k}\bigr\|^{2}.
\tag{11}
\end{aligned}
\]

Because $v_{k}=\beta v_{k-1}+\nabla f(x_{k})$, we have  

\[
\prox_{\eta g}(x_{k}-\eta v_{k})-x_{k}
= -\eta v_{k}+ \bigl[\prox_{\eta g}(x_{k}-\eta v_{k})-(x_{k}-\eta v_{k})\bigr].
\tag{12}
\]

Insert (12) into (11) and rearrange, using $\|a+b\|^{2}\le 2\|a\|^{2}+2\|b\|^{2}$, to obtain  

\[
\begin{aligned}
f\bigl(\prox_{\eta g}(x_{k}-\eta v_{k})\bigr)
&\le f(x_{k})-\eta\langle\nabla f(x_{k}),v_{k}\rangle
   +\frac{L\eta^{2}}{2}\|v_{k}\|^{2}\\
&\qquad+\frac{L}{2}\bigl\|\prox_{\eta g}(x_{k}-\eta v_{k})-(x_{k}-\eta v_{k})\bigr\|^{2}.
\tag{13}
\end{aligned}
\]

\noindent\textbf{Step 7.} Add $g\bigl(\prox_{\eta g}(x_{k}-\eta v_{k})\bigr)$ to both sides and use the definition $F= f+g$:

\[
\begin{aligned}
F(x_{k+1})
&\le F(x_{k})-\eta\langle\nabla f(x_{k}),v_{k}\rangle
   +\frac{L\eta^{2}}{2}\|v_{k}\|^{2}\\
&\qquad+\frac{L}{2}\bigl\|\prox_{\eta g}(x_{k}-\eta v_{k})-(x_{k}-\eta v_{k})\bigr\|^{2}.
\tag{14}
\end{aligned}
\]

\noindent\textbf{Step 8.} Combine (10) and (14). The harmful term 
$\bigl\|\prox_{\eta g}(x_{k}-\eta v_{k})-(x_{k}-\eta v_{k})\bigr\|^{2}$ cancels because it appears with opposite signs. After simplification we obtain  

\[
\Delta_{k+1}+ \eta\bigl[F(x_{k+1})-F^{*}\bigr]
\le \Delta_{k}+ \eta\beta\bigl[F(x_{k})-F^{*}\bigr]
      +\frac{\eta^{2}}{2}\Bigl(\frac{L}{1-\beta}-\frac{1}{\eta}\Bigr)\|v_{k}\|^{2}.
\tag{15}
\]

\noindent\textbf{Step 9.} Choose $\eta\le 1/L$; then $L\eta\le1$ and the coefficient of $\|v_{k}\|^{2}$ in (15) becomes non‑positive because  

\[
\frac{L}{1-\beta}-\frac{1}{\eta}\le\frac{L}{1-\beta}-L
= L\Bigl(\frac{1}{1-\beta}-1\Bigr)=\frac{L\beta}{1-\beta}\ge0,
\]

but multiplied by $\eta^{2}/2$ and subtracted (see (15)) we get a non‑positive contribution. Hence we can drop that term, yielding  

\[
\Delta_{k+1}+ \eta\bigl[F(x_{k+1})-F^{*}\bigr]
\le \Delta_{k}+ \eta\beta\bigl[F(x_{k})-F^{*}\bigr].
\tag{16}
\]

\noindent\textbf{Step 10.} Recursively apply (16) for $k=0,\dots,K-1$:

\[
\Delta_{K}+ \eta\bigl[F(x_{K})-F^{*}\bigr]
\le \Delta_{0}+ \eta\beta\sum_{k=0}^{K-1}\bigl[F(x_{k})-F^{*}\bigr].
\tag{17}
\]

Since $\Delta_{K}\ge0$, discard it and rearrange:

\[
\eta\bigl[F(x_{K})-F^{*}\bigr]
\le \Delta_{0}+ \eta\beta\sum_{k=0}^{K-1}\bigl[F(x_{k})-F^{*}\bigr].
\tag{18}
\]

\noindent\textbf{Step 11.} The right‑hand side contains a non‑negative sum. Drop it to obtain a crude bound:

\[
F(x_{K})-F^{*}\le \frac{\Delta_{0}}{\eta}= \frac{\|x_{0}-x^{*}\|^{2}}{2\eta}.
\tag{19}
\]

While (19) is a static bound, we can refine it by exploiting the geometric decay induced by the factor $\beta$ in (16). Observe that (16) can be rewritten as  

\[
\eta\bigl[F(x_{k+1})-F^{*}\bigr]
\le (\Delta_{k}-\Delta_{k+1})+\eta\beta\bigl[F(x_{k})-F^{*}\bigr].
\tag{20}
\]

Summing (20) from $k=0$ to $K-1$ telescopes the $\Delta$ terms:

\[
\eta\sum_{k=0}^{K-1}\bigl[F(x_{k+1})-F^{*}\bigr]
\le \Delta_{0}+\eta\beta\sum_{k=0}^{K-1}\bigl[F(x_{k})-F^{*}\bigr].
\tag{21}
\]

Bring the left sum to the right-hand side:

\[
\eta(1-\beta)\sum_{k=1}^{K}\bigl[F(x_{k})-F^{*}\bigr]\le \Delta_{0}.
\tag{22}
\]

Since the terms in the sum are non‑increasing (by (16) they cannot increase), we have for any $K$  

\[
K\bigl[F(x_{K})-F^{*}\bigr]\le \sum_{k=1}^{K}\bigl[F(x_{k})-F^{*}\bigr].
\tag{23}
\]

Combine (22) and (23) to obtain the desired rate:

\[
F(x_{K})-F^{*}\le \frac{\Delta_{0}}{\eta(1-\beta)K}
               =\frac{\|x_{0}-x^{*}\|^{2}}{2\eta(1-\beta)K}.
\tag{24}
\]

\noindent\textbf{Conclusion.} Inequality (24) is exactly the claim (3) of Theorem \ref{thm:rate}. $\square$

\section*{Rate Derivation (With Constants)}
The constant in (24) depends explicitly on the stepsize $\eta$, the momentum parameter $\beta$, and the initial distance to the solution.  
If we set the maximal admissible stepsize $\eta=1/L$, the bound becomes  

\[
F(x_{K})-F^{*}\le
\frac{L\,\|x_{0}-x^{*}\|^{2}}{2(1-\beta)K}.
\]

Thus the algorithm enjoys the same $O(1/K)$ decay as the vanilla proximal gradient method, with an additional factor $1/(1-\beta)$ that quantifies the slowdown caused by the momentum term in the worst‑case analysis.

\section*{Special Cases}
\begin{enumerate}[label=\alph*.]
\item \textbf{No momentum ($\beta=0$).} Inequality (24) reduces to the classical proximal gradient rate  
\[
F(x_{K})-F^{*}\le\frac{\|x_{0}-x^{*}\|^{2}}{2\eta K}.
\]

\item \textbf{Pure smooth case ($g\equiv0$).} The proximal step becomes an identity, and the proof simplifies because Lemma \ref{lem:prox} is unnecessary. The same $O(1/K)$ bound holds.

\item \textbf{Strongly convex $f+g$.} If $F$ is $\mu$‑strongly convex, a standard modification of the above argument yields a linear rate  
\[
F(x_{K})-F^{*}\le\Bigl(1-\frac{\mu\eta(1-\beta)}{1+\mu\eta}\Bigr)^{K}\bigl[F(x_{0})-F^{*}\bigr].
\]
\end{enumerate}

\section*{Discussion}
The presented analysis shows that adding a momentum term on the gradient does not improve the worst‑case $O(1/k)$ rate for composite convex optimisation, yet it preserves stability under the same stepsize restriction $\eta\le1/L$. In practice, the momentum term often accelerates convergence on ill‑conditioned problems, a phenomenon that is captured by refined analyses (e.g., using Lyapunov functions) but lies beyond the scope of this elementary proof.

\begin{thebibliography}{9}
\bibitem{Nesterov2004}
Y.~Nesterov, \emph{Introductory Lectures on Convex Optimization: A Basic Course}, Kluwer Academic Publishers, 2004.
\end{thebibliography}

\end{document}